\documentclass[12pt]{article}
\usepackage[margin=1in]{geometry}
\usepackage{kotex}
\usepackage{setspace} 
\doublespacing

\title{마지막 남은 비대칭: 공정성, 경쟁, 그리고 출산의 미래}
\author{임강빈\thanks{Email: yim@sch.ac.kr (순천향대학교 교수)}}
\date{2026년 05월}

\begin{document}
\maketitle

수십 년 동안 출산율 하락은 거의 전적으로 경제의 언어로 설명되어 왔다. 집값은 너무 비싸고, 교육비는 과중하며, 고용은 불안정하고, 아이를 기르는 데 드는 기회비용은 계속 커진다. 각 국가의 정부들은 이에 발맞춰 보조금과 세제 혜택, 보육 프로그램, 그리고 점점 확대되는 가족 지원 체계로 대응해 왔다. 그러나 세계에서 가장 선진화된 많은 사회에서 출생률은 여전히 떨어지고 있다. 추세를 되돌리기 위해 쏟아부은 그 많은 재정에도 불구하고 하락이 멈추지 않는다는 사실은, 다루기 까다로운 한 가지 가능성을 떠올리게 한다. 어쩌면 이 문제는 순수하게 경제적인 것이 아닐지도 모른다는 것, 그리고 비용이라는 익숙한 표면 아래에서 더 깊은 무언가가 작동하고 있을지도 모른다는 것이다.

출산에 관한 대부분의 논의는 비용에 집중한다. 문명 자체의 구조를 들여다보는 경우는 훨씬 드물다. 그러나 문명은 놀라울 만큼 일관된 방향으로 움직여 왔다. 그것은 꾸준히 비대칭을 줄여 왔다는 것이다. 지배자와 피지배자 사이의 거리는 좁혀졌다. 귀족과 평민을 가르던 선은 흐려졌다. 교육은 더 널리 열렸고, 정치 참여는 더 포용적이 되었으며, 경제적 삶은 물려받은 특권에 덜 의존하게 되었다. 이렇게 보면 인류 역사의 상당 부분은 공정성이 확장되어 온 하나의 긴 과정으로 읽을 수 있다.

그러나 공정성에는 쉽게 간과되는 결과가 따른다. 공정성은 경쟁을 없애는 것이 아니라 보편화한다. 엄격하게 계층화된 사회에서는 대다수가 경쟁이 시작되기도 전에 배제된다. 더 대칭적인 사회에서는 훨씬 많은 사람이 같은 무대에 들어선다. 그 결과는 더 적은 경쟁이 아니라 더 많은 경쟁이다. 더 넓고, 더 빽빽하며, 더 공유된 경쟁이다. 따라서 현대 세계는 단지 이전보다 더 공정한 세계인 것만이 아니다. 바로 그 공정함 때문에 더 경쟁적인 세계이기도 하다.

이 점이 가장 분명하게 드러나는 곳은 남성과 여성의 관계가 변해 온 양상이다. 인류 역사의 대부분에서 두 성은 상당히 다른 사회적 역할을 맡았고, 그 이유는 단지 문화적인 것만이 아니라 기능적인 것이기도 했다. 생존은 신체적 힘, 환경적 위험에의 노출, 전쟁과 사냥, 그리고 노동집약적 생산에 크게 의존했다. 남성과 여성이 서로 다른 경로로 사회에 들어선 것은, 사회 자체가 서로 다른 요구를 중심으로 조직되어 있었기 때문이다. 그 배치가 정당했는가는 별개의 문제다. 여기서 중요한 것은 경쟁이 대체로 분리되어 있었다는 사실이다. 남성과 여성은 서로 맞붙기보다 흔히 서로 다른 영역 안에서 겨루었다.

현대 문명은 그 구분을 하나씩 허물어 왔다. 산업화는 신체적 힘의 가치를 떨어뜨렸다. 대중 교육은 지식에 대한 접근을 넓혔다. 디지털 기술은 정보의 장벽을 낮췄다. 직업적 기회와 경제적 독립, 정치적 권리가 모두 확장되었다. 그에 따라 남성과 여성은 점점 같은 경쟁 체계 안에서 함께 활동하게 되었다. 그들은 같은 학교에 다니고, 같은 대학 입학 자리를 두고, 같은 일자리와 같은 승진과 같은 지도적 위치를, 그리고 같은 성공의 정의를 두고 경쟁한다. 방향은 분명하다. 현대 사회는 더 큰 평등을 만들어 낼 뿐 아니라, 하나의 공유된 경쟁의 무대를 짜 올리고 있다.

인공지능과 로봇공학은 이제 이 수렴을 가속하고 있다. 기계는 한때 남성의 힘에 유리했던 물리적 노동의 상당 부분을 이미 흡수했고, 알고리즘은 또 다른 형태의 차이를 만들어 내던 인지적 노동을 자동화하기 시작했다. 로봇은 신체적으로 고된 일을 떠맡고, 소프트웨어는 분석적인 일을 떠맡는다. 이 기술들이 발전할수록 생물학적 성은 경제적 생산성과 직업적 성취에 점점 덜 중요해진다. 장기적인 함의는 충격적이다. 인류 역사상 처음으로, 남성과 여성은 거의 완전히 대칭적인 사회 구조 안에서 경쟁하게 될지도 모른다.

그리고 바로 이 지점, 거의 완전한 대칭의 문턱에서 하나의 역설이 떠오른다. 비대칭이 하나씩 사라질수록, 단 하나의 비대칭이 더욱 또렷해진다. 바로 출산이다. 임신은 여성의 몸 안에서 일어나고, 출산 또한 그러하다. 어떤 정치적 개혁도 이를 지울 수 없고, 어떤 경제 정책도 이를 재분배할 수 없으며, 어떤 법적 틀도 이를 완전히 평등하게 만들 수 없다. 그것은 생물학적 사실로 남는다. 그리고 역설적이게도, 사회가 대칭적이 될수록 그 사실은 더욱 선명하게 두드러진다.

이전 사회에서는 비대칭이 도처에 있었다. 남성은 전쟁과 사냥, 위험한 노동을 통해 더 큰 신체적 위험을 짊어지는 경우가 많았고, 여성은 임신과 출산, 양육의 부담을 졌다. 사회가 동시에 수많은 비대칭 위에 세워져 있었기에, 어느 하나가 집단적 의식을 지배하지는 않았다. 오늘날 그 비대칭들은 대부분 약해졌다. 교육적·정치적·직업적·경제적 비대칭이 모두 그러하다. 그것들이 물러나면서 출산은 점점 홀로 남겨진다. 체계에서 제거될 수 없는 몇 안 되는 비대칭 가운데 하나로서.

이것의 의미는 생물학을 넘어선다. 공정한 경쟁을 중심으로 조직된 사회에서, 출산은 그 외의 모든 면에서 대칭적인 무대 안의 비대칭적 부담으로 경험될 수 있다. 이는 여성이 출산을 불공정하다고 판단하여 의식적으로 거부한다는 뜻이 아니다. 인간의 결정은 그보다 훨씬 더 얽혀 있다. 경제적 조건이 작용하고, 문화적 가치와 개인의 열망, 공공 정책 또한 작용한다. 주장은 더 미묘하다. 확장되는 공정성과 격화되는 경쟁이라는 배경 앞에 놓일 때, 출산의 의미가 달라질 수 있다는 것이다. 한때 수많은 비대칭 가운데 하나였던 것이 이제 몇 안 되는 비대칭 가운데 하나로 서 있고, 한때 분리된 사회적 경로 안에 있던 것이 이제 공유된 경로 안에 놓여 있다. 두 가지 이유 모두에서 그것은 더 잘 보이게 된다. 그리고 더 잘 보이게 된 것은 더 중대해질 수도 있다.

이는 출산율 하락이 가장 발전한 사회들에서 가장 심각하게 나타나는 까닭을 설명하는 데 도움이 될지도 모른다. 그런 사회는 단지 더 부유한 것이 아니라 더 대칭적이고 더 경쟁적이다. 그들은 더 큰 교육 기회와 더 큰 직업적 이동성, 더 큰 개인적 자율성을 제공한다. 그리고 점점 더, 남성과 여성을 같은 평가 체계 안에 놓는다. 자유와 공정을 넓히는 바로 그 힘이, 동시에 생물학적 비대칭에 대한 자각을 날카롭게 벼릴 수 있다.

이 자리에서 보면, 출산율 하락은 단순한 인구학적 사건이라기보다 두 역사적 궤적 사이의 긴장처럼 보인다. 공정과 더 넓은 경쟁을 향한 문명의 오랜 운동, 그리고 인간 재생산의 생물학적 구조가 그것이다. 역사의 대부분 동안 이 둘은 사회가 품고 있던 수많은 비대칭에 완충되어 조용히 공존했다. 문명이 그 비대칭들을 제거해 갈수록 재생산 자체가 점점 더 예외적인 일이 된다.

따라서 발전한 사회가 마주한 질문은 정책 입안자들이 흔히 던지는 질문과 다를지도 모른다. 아이를 갖는 경제적 비용을 어떻게 낮출 것인가를 묻는 것으로는 충분하지 않을 수 있다. 더 깊은 질문은 이런 것일 테다. 공정과 보편적 경쟁 위에 세워진 사회가, 완전히 제거할 수 없는 생물학적 비대칭과 마주할 때 무슨 일이 일어나는가? 인공지능과 로봇공학은 이 물음을 더욱 날카롭게 만들 뿐이다. 기술이 노동과 역량, 사회적 참여에서의 구분을 계속 녹여 갈수록, 생물학적 비대칭의 가시성은 더 높아질 수 있다. 결국 출산의 미래는 경제만의 문제라기보다, 문명이 자신의 마지막 남은 비대칭에 어떻게 응답하기로 하는가의 문제일지도 모른다.

\end{document}
